%=========================================================
% APÉNDICE: INTEGRALES Y FUNCIONES TRIGONOMÉTRICAS INVERSAS
% Material para incorporar a las notas de Cálculo Integral
% No utiliza derivación implícita.
% Requiere: amsmath, amssymb, amsthm, mathtools, enumitem, booktabs
% y los entornos definicion, ejemplo, ejercicio, observacion, advertencia
%=========================================================

\appendix
\section{Integrales que producen funciones trigonométricas inversas}
\label{ap:integrales-trig-inversas}

\subsection{Propósito del apéndice}

En algunas integrales aparecen expresiones que no conducen a potencias,
logaritmos o funciones trigonométricas ordinarias, sino a
\emph{funciones trigonométricas inversas}. Las dos formas que aparecen con
mayor frecuencia en un primer curso son
\[
\frac{1}{\sqrt{1-x^2}}
\qquad\text{y}\qquad
\frac{1}{1+x^2}.
\]

El objetivo de este apéndice no es agregar fórmulas aisladas a una tabla de
integración, sino aprender a reconocer su estructura, reconstruir sus formas
escaladas y decidir cuándo es necesario realizar primero una transformación
algebraica o una sustitución.

\begin{observacion}
En este apéndice no se utiliza derivación implícita. Se parte de las fórmulas
de derivación de las funciones trigonométricas inversas, previamente conocidas,
y se emplean lectura inversa de la derivada, álgebra, regla de la cadena y
sustitución.
\end{observacion}

%---------------------------------------------------------
\subsection{Punto de partida: leer una derivada en sentido inverso}
%---------------------------------------------------------

Recordemos las dos derivadas fundamentales
\[
\frac{d}{dx}\arcsin x=\frac{1}{\sqrt{1-x^2}},
\qquad -1<x<1,
\]
y
\[
\frac{d}{dx}\arctan x=\frac{1}{1+x^2},
\qquad x\in\mathbb{R}.
\]

Toda fórmula de derivación puede leerse en sentido inverso como una fórmula
de integración. Por tanto,
\[
\boxed{\int\frac{dx}{\sqrt{1-x^2}}=\arcsin x+C}
\]
y
\[
\boxed{\int\frac{dx}{1+x^2}=\arctan x+C.}
\]

Éstas son las dos parejas que conviene dominar. Las formas más generales se
pueden reconstruir a partir de ellas.

%---------------------------------------------------------
\subsection{Forma de arcoseno: \texorpdfstring{$\sqrt{a^2-x^2}$}{sqrt(a2-x2)}}
%---------------------------------------------------------

Queremos reconstruir
\[
I=\int\frac{dx}{\sqrt{a^2-x^2}},\qquad a>0.
\]

El primer paso es hacer visible la estructura $1-u^2$. Factorizamos $a^2$
dentro de la raíz:
\begin{align*}
\sqrt{a^2-x^2}
&=\sqrt{a^2\left(1-\frac{x^2}{a^2}\right)}\\
&=a\sqrt{1-\left(\frac{x}{a}\right)^2},
\end{align*}
porque $a>0$. Entonces
\[
I=\frac1a\int
\frac{dx}{\sqrt{1-\left(\frac{x}{a}\right)^2}}.
\]

Tomamos
\[
u=\frac{x}{a},\qquad du=\frac1a\,dx.
\]
Por tanto,
\[
I=\int\frac{du}{\sqrt{1-u^2}}
=\arcsin u+C.
\]
Regresando a $x$,
\[
\boxed{\int\frac{dx}{\sqrt{a^2-x^2}}
=\arcsin\left(\frac{x}{a}\right)+C,\qquad a>0.}
\]

\begin{observacion}
En esta fórmula no aparece un factor exterior $1/a$. El factor $1/a$ que
surge al factorizar la raíz coincide exactamente con el que aparece en
$du=(1/a)dx$.
\end{observacion}

\subsubsection*{Ejemplo 1: reconocimiento directo}
Calculemos
\[
\int\frac{dx}{\sqrt{25-x^2}}.
\]
Como $25=5^2$,
\[
\int\frac{dx}{\sqrt{25-x^2}}
=\boxed{\arcsin\left(\frac{x}{5}\right)+C}.
\]

\subsubsection*{Ejemplo 2: primero hay que reescribir}
Consideremos
\[
\int\frac{dx}{\sqrt{16-9x^2}}.
\]
La expresión todavía no tiene exactamente la forma $a^2-x^2$. Escribimos
\[
16-9x^2=16\left(1-\left(\frac{3x}{4}\right)^2\right).
\]
Así,
\[
\sqrt{16-9x^2}=4\sqrt{1-\left(\frac{3x}{4}\right)^2}
\]
y
\[
I=\frac14\int\frac{dx}{\sqrt{1-(3x/4)^2}}.
\]
Tomamos
\[
u=\frac{3x}{4},\qquad du=\frac34\,dx,
\qquad dx=\frac43\,du.
\]
Entonces
\begin{align*}
I
&=\frac14\frac43\int\frac{du}{\sqrt{1-u^2}}\\
&=\frac13\arcsin u+C.
\end{align*}
Por tanto,
\[
\boxed{\int\frac{dx}{\sqrt{16-9x^2}}
=\frac13\arcsin\left(\frac{3x}{4}\right)+C.}
\]

\subsubsection*{Ejemplo 3: una constante adicional en el numerador}
\[
\int\frac{6\,dx}{\sqrt{49-4x^2}}.
\]
Como
\[
49-4x^2=49\left(1-\left(\frac{2x}{7}\right)^2\right),
\]
tenemos
\[
I=\frac67\int\frac{dx}{\sqrt{1-(2x/7)^2}}.
\]
Con
\[
u=\frac{2x}{7},\qquad dx=\frac72\,du,
\]
resulta
\[
I=3\int\frac{du}{\sqrt{1-u^2}}
=\boxed{3\arcsin\left(\frac{2x}{7}\right)+C}.
\]

%---------------------------------------------------------
\subsection{Forma de arcotangente: \texorpdfstring{$a^2+x^2$}{a2+x2}}
%---------------------------------------------------------

Partimos de
\[
I=\int\frac{dx}{a^2+x^2},\qquad a>0.
\]
Factorizamos $a^2$:
\[
a^2+x^2=a^2\left(1+\left(\frac{x}{a}\right)^2\right).
\]
Entonces
\[
I=\frac1{a^2}\int
\frac{dx}{1+(x/a)^2}.
\]
Tomamos
\[
u=\frac{x}{a},\qquad dx=a\,du.
\]
Por tanto,
\begin{align*}
I
&=\frac1{a^2}\int\frac{a\,du}{1+u^2}\\
&=\frac1a\arctan u+C.
\end{align*}
Así,
\[
\boxed{\int\frac{dx}{a^2+x^2}
=\frac1a\arctan\left(\frac{x}{a}\right)+C,\qquad a>0.}
\]

\begin{observacion}
Aquí sí permanece un factor $1/a$. Ésta es una diferencia importante entre
las formas escaladas del arcoseno y de la arcotangente.
\end{observacion}

\subsubsection*{Ejemplo 4: reconocimiento directo}
\[
\int\frac{dx}{9+x^2}
=\boxed{\frac13\arctan\left(\frac{x}{3}\right)+C}.
\]

\subsubsection*{Ejemplo 5: coeficiente en \texorpdfstring{$x^2$}{x2}}
Consideremos
\[
\int\frac{dx}{4+25x^2}.
\]
Escribimos
\[
4+25x^2=4\left(1+\left(\frac{5x}{2}\right)^2\right).
\]
Así,
\[
I=\frac14\int\frac{dx}{1+(5x/2)^2}.
\]
Sea
\[
u=\frac{5x}{2},\qquad dx=\frac25\,du.
\]
Entonces
\[
I=\frac1{10}\int\frac{du}{1+u^2}
=\boxed{\frac1{10}\arctan\left(\frac{5x}{2}\right)+C}.
\]

\subsubsection*{Ejemplo 6: completar el cuadrado}
Calculemos
\[
\int\frac{dx}{x^2+6x+13}.
\]
La estructura de arcotangente no es visible de inmediato. Completamos el
cuadrado:
\begin{align*}
x^2+6x+13
&=x^2+6x+9+4\\
&=(x+3)^2+2^2.
\end{align*}
Por tanto,
\[
I=\int\frac{dx}{(x+3)^2+2^2}.
\]
Tomamos $u=x+3$, $du=dx$:
\[
I=\int\frac{du}{u^2+2^2}
=\boxed{\frac12\arctan\left(\frac{x+3}{2}\right)+C}.
\]

\subsubsection*{Ejemplo 7: completar el cuadrado con coeficiente principal}
Consideremos
\[
\int\frac{dx}{2x^2-8x+10}.
\]
Factorizamos primero el $2$:
\[
2x^2-8x+10=2(x^2-4x+5).
\]
Completamos el cuadrado:
\[
x^2-4x+5=(x-2)^2+1.
\]
Entonces
\begin{align*}
I
&=\frac12\int\frac{dx}{(x-2)^2+1}\\
&=\boxed{\frac12\arctan(x-2)+C}.
\end{align*}

%---------------------------------------------------------
\subsection{Arcocoseno: una forma equivalente}
%---------------------------------------------------------

También sabemos que
\[
\frac{d}{dx}\arccos x=-\frac1{\sqrt{1-x^2}}.
\]
Por tanto,
\[
\int\frac{-dx}{\sqrt{1-x^2}}=\arccos x+C,
\]
o equivalentemente,
\[
\int\frac{dx}{\sqrt{1-x^2}}=-\arccos x+C.
\]
Esto no contradice la fórmula con arcoseno, porque
\[
\arcsin x+\arccos x=\frac\pi2.
\]
En consecuencia, $\arcsin x$ y $-\arccos x$ difieren solamente en una
constante.

\begin{observacion}
Para integración práctica utilizaremos preferentemente la forma con
$\arcsin$, pues permite mantener una sola fórmula de referencia para la
estructura $1/\sqrt{1-x^2}$.
\end{observacion}

%---------------------------------------------------------
\subsection{Arcosecante y la importancia del valor absoluto}
%---------------------------------------------------------

Otra fórmula que aparece en tablas de integración es
\[
\frac{d}{dx}\operatorname{arcsec}x
=\frac1{|x|\sqrt{x^2-1}},\qquad |x|>1.
\]
Leyéndola en sentido inverso,
\[
\boxed{\int\frac{dx}{|x|\sqrt{x^2-1}}
=\operatorname{arcsec}x+C.}
\]
Una forma escalada frecuente es
\[
\boxed{\int\frac{dx}{|x|\sqrt{x^2-a^2}}
=\frac1a\operatorname{arcsec}\left(\frac{|x|}{a}\right)+C,\qquad a>0,}
\]
con las consideraciones correspondientes de dominio y de la convención usada
para $\operatorname{arcsec}$.

\begin{advertencia}
El valor absoluto no es un adorno. Aparece porque la derivada de la función
arcosecante debe tratar correctamente las dos regiones $x>1$ y $x<-1$. En un
primer curso conviene dar mayor prioridad operativa a las formas de arcoseno
y arcotangente.
\end{advertencia}

%---------------------------------------------------------
\subsection{Reconocimiento estructural: no basta con mirar la fórmula}
%---------------------------------------------------------

Las formas anteriores se vuelven realmente útiles cuando aparecen
composiciones. Por la regla de la cadena, las estructuras que debemos buscar
son
\[
\boxed{\frac{f'(x)}{\sqrt{1-f(x)^2}}}
\quad\longrightarrow\quad
\arcsin(f(x))+C,
\]
y
\[
\boxed{\frac{f'(x)}{1+f(x)^2}}
\quad\longrightarrow\quad
\arctan(f(x))+C.
\]

La pregunta correcta no es solamente ``¿veo un $1-x^2$?'' sino
\begin{quote}
¿puedo identificar una expresión interior $f(x)$ y aparece también, salvo una
constante, su derivada $f'(x)$?
\end{quote}

\subsubsection*{Ejemplo 8: composición que produce arcoseno}
\[
I=\int\frac{2x\,dx}{\sqrt{1-x^4}}.
\]
Observamos que $x^4=(x^2)^2$. Tomamos
\[
u=x^2,\qquad du=2x\,dx.
\]
Entonces
\[
I=\int\frac{du}{\sqrt{1-u^2}}
=\boxed{\arcsin(x^2)+C}.
\]

\subsubsection*{Ejemplo 9: composición que produce arcotangente}
\[
I=\int\frac{3x^2\,dx}{1+x^6}.
\]
Como $x^6=(x^3)^2$, tomamos
\[
u=x^3,\qquad du=3x^2\,dx.
\]
Así,
\[
I=\int\frac{du}{1+u^2}
=\boxed{\arctan(x^3)+C}.
\]

\subsubsection*{Ejemplo 10: ajustar una constante}
\[
I=\int\frac{x\,dx}{\sqrt{1-4x^4}}.
\]
La expresión interior natural es $u=2x^2$, pues $u^2=4x^4$. Entonces
\[
du=4x\,dx,
\qquad x\,dx=\frac14du.
\]
Por tanto,
\[
I=\frac14\int\frac{du}{\sqrt{1-u^2}}
=\boxed{\frac14\arcsin(2x^2)+C}.
\]

%---------------------------------------------------------
\subsection{Una estrategia de decisión}
%---------------------------------------------------------

Cuando una integral parezca relacionada con funciones trigonométricas
inversas, conviene seguir el siguiente procedimiento.

\begin{enumerate}
\item \textbf{Simplificar algebraicamente} el integrando si es posible.
\item Buscar una estructura del tipo
\[
\sqrt{a^2-u^2}
\qquad\text{o}\qquad
a^2+u^2.
\]
\item Identificar la expresión interior $u=f(x)$.
\item Comprobar si aparece $f'(x)\,dx$, quizá multiplicado por una constante.
\item Si existe un trinomio cuadrático, considerar \textbf{completar el cuadrado}.
\item Realizar la sustitución y reducir a una de las dos formas fundamentales.
\item Regresar a la variable original.
\item \textbf{Verificar derivando} el resultado.
\end{enumerate}

\subsection{Tabla de reconocimiento rápido}

\begin{center}
\renewcommand{\arraystretch}{1.6}
\begin{tabular}{>{\centering\arraybackslash}p{0.29\textwidth}
                >{\centering\arraybackslash}p{0.29\textwidth}
                >{\centering\arraybackslash}p{0.30\textwidth}}
\toprule
\textbf{Estructura} & \textbf{Sustitución o reconocimiento} & \textbf{Resultado típico}\\
\midrule
$\displaystyle \frac{1}{\sqrt{1-u^2}}$
& identificar $u=f(x)$
& $\arcsin u+C$\\
$\displaystyle \frac{1}{1+u^2}$
& identificar $u=f(x)$
& $\arctan u+C$\\
$\displaystyle \frac{1}{\sqrt{a^2-x^2}}$
& $u=x/a$
& $\arcsin(x/a)+C$\\
$\displaystyle \frac{1}{a^2+x^2}$
& $u=x/a$
& $\frac1a\arctan(x/a)+C$\\
$\displaystyle \frac{f'(x)}{\sqrt{1-f(x)^2}}$
& $u=f(x)$
& $\arcsin(f(x))+C$\\
$\displaystyle \frac{f'(x)}{1+f(x)^2}$
& $u=f(x)$
& $\arctan(f(x))+C$\\
\bottomrule
\end{tabular}
\end{center}

%---------------------------------------------------------
\subsection{Errores frecuentes}
%---------------------------------------------------------

\begin{enumerate}
\item Confundir
\[
\int\frac{dx}{a^2+x^2}
\]
con una integral logarítmica. Para obtener $\ln(a^2+x^2)$ tendría que aparecer
en el numerador una constante múltiplo de $2x$.

\item Olvidar el factor $1/a$ en
\[
\int\frac{dx}{a^2+x^2}.
\]

\item Suponer que
\[
\sqrt{a^2-x^2}=a-x.
\]
Esta igualdad es falsa en general.

\item Aplicar una fórmula antes de normalizar la expresión. Por ejemplo,
$\sqrt{16-9x^2}$ no es todavía $\sqrt{a^2-x^2}$ con la variable $x$ sin
modificación: debe hacerse visible $(3x/4)^2$.

\item Completar incorrectamente el cuadrado. Siempre conviene comprobar la
identidad expandiendo el resultado.

\item Hacer una sustitución $u=f(x)$ sin verificar que el diferencial $du$
aparezca en el integrando, al menos salvo un factor constante.
\end{enumerate}

%---------------------------------------------------------
\subsection{Ejercicios propuestos}
%---------------------------------------------------------

\subsubsection*{Nivel 1: reconocimiento directo}
Calcule:
\begin{enumerate}[label=\arabic*.]
\item $\displaystyle \int\frac{dx}{\sqrt{36-x^2}}$.
\item $\displaystyle \int\frac{dx}{16+x^2}$.
\item $\displaystyle \int\frac{dx}{\sqrt{9-x^2}}$.
\item $\displaystyle \int\frac{dx}{25+x^2}$.
\end{enumerate}

\subsubsection*{Nivel 2: normalización y constantes}
\begin{enumerate}[label=\arabic*.,resume]
\item $\displaystyle \int\frac{dx}{\sqrt{25-4x^2}}$.
\item $\displaystyle \int\frac{dx}{9+16x^2}$.
\item $\displaystyle \int\frac{5\,dx}{\sqrt{49-25x^2}}$.
\item $\displaystyle \int\frac{6\,dx}{4+9x^2}$.
\end{enumerate}

\subsubsection*{Nivel 3: composición}
\begin{enumerate}[label=\arabic*.,resume]
\item $\displaystyle \int\frac{2x\,dx}{1+x^4}$.
\item $\displaystyle \int\frac{3x^2\,dx}{\sqrt{1-x^6}}$.
\item $\displaystyle \int\frac{e^x\,dx}{1+e^{2x}}$.
\item $\displaystyle \int\frac{\cos x\,dx}{\sqrt{1-\sin^2x}}$, indicando en qué intervalos se simplifica sin ambigüedad.
\end{enumerate}

\subsubsection*{Nivel 4: completar el cuadrado}
\begin{enumerate}[label=\arabic*.,resume]
\item $\displaystyle \int\frac{dx}{x^2+4x+8}$.
\item $\displaystyle \int\frac{dx}{x^2-6x+18}$.
\item $\displaystyle \int\frac{dx}{2x^2+8x+10}$.
\item $\displaystyle \int\frac{dx}{4x^2-8x+8}$.
\end{enumerate}

%---------------------------------------------------------
\subsection{Respuestas de los ejercicios}
%---------------------------------------------------------

\begin{enumerate}[label=\arabic*.]
\item $\displaystyle \arcsin(x/6)+C$.
\item $\displaystyle \frac14\arctan(x/4)+C$.
\item $\displaystyle \arcsin(x/3)+C$.
\item $\displaystyle \frac15\arctan(x/5)+C$.
\item $\displaystyle \frac12\arcsin(2x/5)+C$.
\item $\displaystyle \frac1{12}\arctan(4x/3)+C$.
\item $\displaystyle \arcsin(5x/7)+C$.
\item $\displaystyle \arctan(3x/2)+C$.
\item $\displaystyle \arctan(x^2)+C$.
\item $\displaystyle \arcsin(x^3)+C$.
\item $\displaystyle \arctan(e^x)+C$.
\item Requiere considerar $\sqrt{1-\sin^2x}=|\cos x|$; la respuesta depende del intervalo considerado.
\item $\displaystyle \frac12\arctan\left(\frac{x+2}{2}\right)+C$.
\item $\displaystyle \frac13\arctan\left(\frac{x-3}{3}\right)+C$.
\item $\displaystyle \frac12\arctan(x+2)+C$.
\item $\displaystyle \frac14\arctan(x-1)+C$.
\end{enumerate}

%---------------------------------------------------------
\subsection{Cierre conceptual}
%---------------------------------------------------------

Las fórmulas de integración asociadas a funciones trigonométricas inversas
son más útiles cuando se entienden como \emph{patrones estructurales}. Las dos
formas fundamentales son
\[
\boxed{\int\frac{du}{\sqrt{1-u^2}}=\arcsin u+C}
\qquad\text{y}\qquad
\boxed{\int\frac{du}{1+u^2}=\arctan u+C.}
\]
A partir de ellas se obtienen las formas escaladas y las composiciones. En
problemas más elaborados, la tarea principal consiste en transformar el
integrando hasta hacer visible una de estas estructuras. Por ello, completar
el cuadrado, factorizar constantes y reconocer una expresión junto con su
derivada son habilidades más importantes que memorizar una lista extensa de
fórmulas.
